% \section{Roadmap of the Thesis}
% \label{sec:intro:roadmap}

% \textbf{\cref{ch:prior}} establishes a conceptual framework for dynamics-constrained motion generation, organizing existing approaches along four design axes---modeling burden, problem dimensionality, system integration, and usability.
% This framework motivates the progression of methods developed throughout the remainder of the thesis.

% \begin{enumerate}
  % \item \textbf{A design framework for dynamics-aware motion generation} (\cref{ch:prior}).
        % A conceptual framework organizing manipulation methods along four design axes---modeling burden, problem dimensionality, system integration, and usability---providing a unified perspective on the tradeoffs underlying dynamics-aware planning systems.

%   \item \textbf{PokeRRT: impulsive non-prehensile manipulation through kinodynamic search} (\cref{ch:poking}).
%         A sampling-based kinodynamic planner that uses physics simulation as a forward model for object rearrangement via impulsive contact, without requiring explicit contact models.
%         Across six scenarios, poking achieves higher success rates and lower task completion times than pushing and grasping baselines.

%   \item \textbf{Kinodynamic planning under coupled robot-object dynamics} (\cref{ch:constrained}).
%         Extensions of kinodynamic planning to manipulation tasks governed by coupled dynamics, including whole-body manipulation under locked multi-joint failures and spill-free liquid transport in cluttered environments.
%         These methods maintain up to 92\% of nominal reachable workspace under failures and expand feasible spill-free solution spaces by up to $3\times$ over prior approaches.
% \end{enumerate}



% \begin{enumerate}
%   \setcounter{enumi}{2}
%   \item \textbf{LazyBoE: multi-query kinodynamic planning via lazy edge bundles} (\cref{ch:multiquery}).
%         A multi-query kinodynamic planner that amortizes dynamics computation across repeated planning problems by precomputing control-parameterized edge bundles over discretized state and control spaces.
%         Lazy propagation and collision checking substantially reduce repeated planning cost while preserving kinodynamic feasibility.
% \end{enumerate}






% \begin{enumerate}
%   \setcounter{enumi}{3}
%   \item \textbf{Motion generation as conditional sampling} (\cref{ch:generation}).
%         A reformulation of dynamics-aware motion generation as conditional trajectory sampling, establishing a conceptual bridge between kinodynamic planning and learned generative models.

%   \item \textbf{Conditional trajectory diffusion for dynamic constraint satisfaction} (\cref{ch:payload}).
%         A diffusion-based trajectory generator conditioned on task-level dynamic properties such as payload mass and embodiment configuration.
%         The method achieves super-nominal payload manipulation---67.6\% workspace accessibility at $3\times$ rated payload capacity---in approximately 10\,ms without runtime constraint checking, and generalizes to fail-active manipulation under actuation failures.

%   \item \textbf{Task-agnostic trajectory seeding for constrained optimization} (\cref{ch:seeding}).
%         A conditioning mechanism that represents kinematic and dynamic constraints as robot state samples in order to generate near-feasible trajectory seeds for constrained optimization.
%         The approach improves solver convergence and feasibility across diverse manipulation tasks without task-specific engineering.
% \end{enumerate}


% \paragraph{Online Constraint Evaluation: Kinodynamic Planning.}
\textbf{Chapters~\ref{ch:poking}--\ref{ch:constrained}} develop methods in which dynamic feasibility is evaluated online: constraint knowledge resides in a physics simulation forward model, and each planning query incurs the full cost of feasibility checking.
\cref{ch:poking} introduces poking as a dynamic non-prehensile manipulation primitive, using impulsive contact within a sampling-based planner to achieve object rearrangement without explicit contact modeling. Across six scenarios, poking achieves higher success rates and lower task completion times than pushing and grasping baselines.
\cref{ch:constrained} extends kinodynamic planning to manipulation tasks governed by coupled robot-object dynamics, including whole-body manipulation under locked multi-joint failures and spill-free liquid transport in cluttered environments. These methods maintain up to 92\% of nominal reachable workspace under failures and expand feasible spill-free solution spaces by up to $3\times$ over prior approaches \ap{fix: move liquid to next section - multiquery}.

% \paragraph{Amortized Feasibility via Reusable Motion Structure.}
\textbf{\cref{ch:multiquery}} introduces the first amortized representation: constraint knowledge is moved upstream into precomputed control-parameterized edge bundles, reducing per-query cost while preserving kinodynamic feasibility through lazy propagation and collision checking.

% \paragraph{Amortized Feasibility via Learned Generative Models.}
\textbf{Chapters~\ref{ch:generation} and \ref{ch:payload}} develop a more aggressive form of amortization: constraint knowledge is absorbed into the parameters of a conditional trajectory distribution during training, eliminating runtime feasibility checking entirely.
\cref{ch:generation} reframes motion generation as conditional sampling, establishing the conceptual bridge between kinodynamic planning and trajectory diffusion models.
\cref{ch:payload} develops a diffusion-based trajectory generator conditioned on task-level dynamic properties such as payload mass, enabling super-nominal payload manipulation without runtime constraint checking. The method achieves super-nominal payload manipulation---67.6\% workspace accessibility at $3\times$ rated payload capacity---in approximately 10\,ms without runtime constraint checking, and generalizes to fail-active manipulation under actuation failures.
A case study within the chapter extends the same conditioning principle to fail-active manipulation under actuation failures, demonstrating adaptation to changes in robot embodiment and dynamic capability.

\textbf{\cref{ch:seeding}} extends the learned representation to constrained trajectory optimization, completing a progression in which the same conditioning principle is applied at different stages of the motion generation pipeline.
A task-agnostic conditioning mechanism represents dynamic constraints as robot state samples to generate near-feasible trajectory seeds that improve optimization convergence and feasibility across diverse manipulation tasks without task-specific engineering.

Finally, \textbf{\cref{ch:discussion}} synthesizes the thesis contributions, discusses their implications for dynamics-aware robot manipulation, and outlines directions for future research.

% \section{Roadmap of the Thesis}
% \label{sec:intro:roadmap}

% \cref{ch:framework} establishes a conceptual framework for dynamics-constrained motion generation, organizing existing approaches along four design axes---modeling burden, problem dimensionality, system integration, and usability.
% This framework motivates the progression of methods developed throughout the remainder of the thesis.

% Chapters~\ref{ch:poking}--\ref{ch:multiquery} investigate dynamics through sampling-based kinodynamic planning.
% These chapters progressively expand the scope of manipulation behaviors while introducing increasing levels of structure reuse and computational amortization.

% \cref{ch:poking} introduces poking as a dynamic non-prehensile manipulation primitive, using impulsive contact within a sampling-based kinodynamic planner to achieve object rearrangement without explicit contact modeling.
% \cref{ch:constrained} extends kinodynamic planning to manipulation tasks governed by coupled robot-object dynamics, including whole-body manipulation under locked multi-joint failures and spill-free liquid transport in cluttered environments.
% \cref{ch:multiquery} then addresses the repeated-query setting through LazyBoE, a multi-query kinodynamic planner that amortizes dynamics computation via precomputed control-parameterized edge bundles combined with lazy propagation and collision checking.

% Chapters~\ref{ch:generation}--\ref{ch:failure} transition from explicit online dynamics reasoning toward learned generative representations of dynamic feasibility.
% \cref{ch:generation} reframes motion generation as conditional sampling, establishing the conceptual bridge between kinodynamic planning and trajectory diffusion models.
% \cref{ch:payload} develops a diffusion-based trajectory generator conditioned on payload mass, enabling super-nominal payload manipulation without runtime constraint checking.
% \cref{ch:failure} generalizes this conditioning principle to embodiment variation through fail-active manipulation under actuation failures.

% \cref{ch:seeding} extends this generative perspective to constrained trajectory optimization by introducing a task-agnostic conditioning mechanism that represents dynamic constraints as robot state samples.
% The resulting method generates near-feasible trajectory seeds that improve optimization convergence and feasibility across diverse manipulation tasks without task-specific engineering.

% Finally, \cref{ch:discussion} synthesizes the thesis contributions, discusses their implications for dynamics-aware robot manipulation, and outlines directions for future research.

% \section{Roadmap of the Thesis}
% \label{sec:intro:roadmap}

% \cref{ch:framework} establishes a conceptual framework for dynamics-constrained motion generation, organizing the landscape of approaches and motivating the progression of methods developed in subsequent chapters.

% Chapters~\ref{ch:poking}--\ref{ch:multiquery} develop sampling-based kinodynamic methods of increasing generality.
% \cref{ch:poking} introduces poking as a dynamic manipulation primitive, using impulsive contact to achieve non-prehensile rearrangement within a sampling-based planner.
% \cref{ch:constrained} extends kinodynamic planning to constrained settings, handling tasks where robot and object dynamics are coupled---such as whole-body contact and liquid-filled container transport.
% \cref{ch:multiquery} addresses the multi-query setting via lazy edge bundles, amortizing planning cost across repeated queries in a shared workspace.

% \cref{ch:generation} marks a transition from sampling to generation, reframing motion planning as conditional sampling and establishing the conceptual bridge to the generative methods that follow.
% \cref{ch:payload} develops diffusion-based trajectory generation conditioned on high-level object properties, demonstrating that dynamic constraints can be implicitly encoded through the conditioning signal; a case study extends this principle to fail-active manipulation under actuation failures.

% \cref{ch:seeding} shifts focus to constrained trajectory optimization, contributing a task-agnostic seeding strategy that substantially improves solver convergence and feasibility without task-specific design.

% \cref{ch:discussion} concludes with a synthesis of the thesis contributions, a discussion of their collective implications for dynamic manipulation, and directions for future work.
